\documentclass[twocolumn, 12pt]{article}

\usepackage[utf8]{inputenc}
\usepackage[english, spanish]{babel}
\usepackage{fullpage}
\usepackage{graphicx}
\usepackage{amsmath}
\usepackage{enumitem}
\usepackage{chngcntr}
\usepackage{setspace}
\usepackage{url}
\usepackage{csquotes}
\usepackage{float}
\usepackage{verbatim}
\usepackage{tabularx}
\usepackage{amsmath}
\usepackage{caption}
\usepackage{bm}
\usepackage{colortbl}
\usepackage{xcolor}

\usepackage{multirow}

% \usepackage{hyperref}

\definecolor{LigthGray}{rgb}{0.7098, 0.7294, 0.7215}
\definecolor{LigthGrayPlus}{rgb}{0.6862, 0.7254, 0.7294}
\definecolor{White}{rgb}{1, 1, 1}
\definecolor{Red}{rgb}{1, 0, 0}

\counterwithin{figure}{section}
\renewcommand{\thesection}{\arabic{section}}
\renewcommand{\thesubsection}{\thesection.\arabic{subsection}}
\renewcommand{\baselinestretch}{1.5}

\usepackage[style=apa, maxnames=6, minnames=3, backend=biber]{biblatex}
\DefineBibliographyStrings{english}{%chktex-file 1 chktex-file 6
      andothers = {\em et\addabbrvspace al\adddot}
}
\addbibresource{./Bibliography/bibliography.bib}

\usepackage{array}

\setlength{\parskip}{0pt}

\raggedbottom{}

\newcommand{\bolditalic}[1]{\textbf{\textit{#1}}}

\begin{document}

\begin{titlepage}
      \centering
      \includegraphics[width=0.3\textwidth]{Images/logo_utb.png}\par\vspace{1cm}
      {\scshape\LARGE Universidad Tecnológica de Bolívar \par}
      \vspace{1cm}

      {\scshape\Large FÍSICA CALOR Y ONDAS \par}
      \vspace{.2cm}

      % chktex-file 8
      {\scshape\Large Grupo 1 \par}
      \vspace{1cm}
      % chktex-file 8
      \slshape {\Large \bfseries{}Informe de Laboratorio No. III\\}
      \slshape {\small \bfseries{}DIFRACCIÓN DE LA LUZ}
      \vspace{2cm}

      \slshape {\itshape{} Mauro González, T00067622 \\}
      \slshape {\itshape{} German De Armas Castaño, T00068765 \\}
      \slshape {\itshape{} Angel Vega Rodriguez, T00068186 \\}
      \slshape {\itshape{} Juan Jose Osorio Ariza, T00067316 \\}
      \slshape {\itshape{} Jorge Alberto Rueda Salgado, T00068722 \\}
      \vfill
      Revisado Por \\
      Duban Andres Paternina Verona\\
      {\large \today\par}
\end{titlepage}

% chktex-file 44
% chktex-file 24

% ! ----------------------------------------------------------------------|>
\section{Introducción}

La difracción de la luz es un fenómeno fascinante que nos
permite comprender cómo la luz se comporta cuando pasa a
través de obstáculos o se dispersa por rendijas estrechas.
Este fenómeno es esencial en la óptica y nos brinda
información valiosa sobre las propiedades de la luz. En
esta experiencia de laboratorio, exploraremos la difracción
de la luz utilizando una rejilla de difracción y un láser
de semiconductor. A través de este experimento, podremos
observar y analizar patrones de interferencia que se forman
cuando la luz se difracta en una rejilla.

% ! ----------------------------------------------------------------------|>
\section{Objetivos}

% + ----------------------------------------|>
\subsection{Objetivo general}

\begin{itemize}[label=$\triangleright$]
      \item El objetivo principal de esta experiencia de laboratorio es
            observar y comprender el fenómeno de difracción de la luz.
            A través de la utilización de una rejilla de difracción y
            un láser de semiconductor, buscaremos analizar y
            caracterizar los patrones de interferencia resultantes de
            la difracción de la luz. Además, se pretende determinar la
            longitud de onda de emisión del láser utilizando los
            patrones de difracción.
\end{itemize}

% + ----------------------------------------|>
\subsection{Objetivos específicos}

\begin{itemize}[label=$\triangleright$]
      \item Utilizar el patrón de difracción para determinar la
            longitud de onda de emisión del láser de semiconductor.

      \item Comparar los resultados experimentales obtenidos en el
            laboratorio con los valores teóricos esperados y verificar
            la precisión del fenómeno de difracción.

      \item Comprender el principio de Huygens y su relación con la
            difracción de la luz.

      \item Explorar la influencia del número de líneas por centímetro
            de la rejilla de difracción en la formación de los patrones
            de interferencia.
\end{itemize}

% ! ----------------------------------------------------------------------|>
\section{Marco Teórico}

 {\Large
  \begin{equation}
        \lambda_{n} = \frac{d \cdot \sen ( \tan^{-1} (\frac{X_{n}}{D}))}{n}
        \label{eq:lambda-n}
  \end{equation}
 }

 {\Large
  \begin{equation}
        \mathcal{G}_x = \frac{\lambda D}{x}
        \label{eq:grosor-cabello}
  \end{equation}
 }

% ! ----------------------------------------------------------------------|>
\section{Montaje Experimental}

% ! ----------------------------------------------------------------------|>
\section{Datos Experimentales}

\begin{table}[H]
      \begin{center}
            \begin{tabularx}{\linewidth}{|>{\centering\arraybackslash}X|>{\centering\arraybackslash}X|>{\centering\arraybackslash}X|}
                  \multicolumn{3}{c}{$D_{1} = 0,7111$ \bolditalic{(m)}}                                                                                \\\hline
                  \multirow{3}{*}{\textbf{n}} & \multirow{3}{*}{$x$ \bolditalic{(m)}} & {\Large $\lambda_{n}$} \bolditalic{(mm)} \break{} Experimental \\\hline
                  1                           & $4,50 \cdot 10^{-02}$                 & $6,32 \cdot 10^{-04}$                                          \\\hline

                  2                           & $9,00 \cdot 10^{-02}$                 & $6,29 \cdot 10^{-04}$                                          \\\hline

                  3                           & $1,37 \cdot 10^{-01}$                 & $6,34 \cdot 10^{-04}$                                          \\\hline

            \end{tabularx}

            \vspace{-.02cm}

            \begin{tabularx}{\linewidth}{|>{\centering\arraybackslash}X|>{\centering\arraybackslash}X|}
                  {\large $\lambda_{Promedio}$} \bolditalic{(mm)} & $6,32 \cdot 10^{-4}$ \\\hline
            \end{tabularx}
      \end{center}
      \caption{}
      \label{tab:d1}
\end{table}

\begin{table}[H]
      \begin{center}
            \begin{tabularx}{\linewidth}{|>{\centering\arraybackslash}X|>{\centering\arraybackslash}X|>{\centering\arraybackslash}X|}
                  \multicolumn{3}{c}{$D_{2} = 0,4500$ \bolditalic{(m)}}                                                                                \\\hline
                  \multirow{3}{*}{\textbf{n}} & \multirow{3}{*}{$x$ \bolditalic{(m)}} & {\Large $\lambda_{n}$} \bolditalic{(mm)} \break{} Experimental \\\hline
                  1                           & $2,80 \cdot 10^{-02}$                 & $6,21 \cdot 10^{-04}$                                          \\\hline

                  2                           & $5,70 \cdot 10^{-02}$                 & $6,30 \cdot 10^{-04}$                                          \\\hline

                  3                           & $8,70 \cdot 10^{-02}$                 & $6,36 \cdot 10^{-04}$                                          \\\hline

            \end{tabularx}

            \vspace{-.02cm}

            \begin{tabularx}{\linewidth}{|>{\centering\arraybackslash}X|>{\centering\arraybackslash}X|}
                  {\large $\lambda_{Promedio}$} \bolditalic{(mm)} & $6,29 \cdot 10^{-4}$ \\\hline
            \end{tabularx}
      \end{center}
      \caption{}
      \label{tab:d2}
\end{table}

\begin{table}[H]
      \begin{center}
            \begin{tabularx}{\linewidth}{|>{\centering\arraybackslash}X|>{\centering\arraybackslash}X|}
                  \multicolumn{2}{c}{Datos (Grosor de un cabello)}                                 \\\hline
                  \textbf{D} \bolditalic{(m)}                                            & $1,48$  \\\hline
                  $\lambda$ \bolditalic{(nm)}~\footnote{Proporcionado por el fabricante} & $632,8$ \\\hline
                  $x$ \bolditalic{(cm)}                                                  & $0,8$   \\\hline
            \end{tabularx}

            \vspace{-.02cm}

            \begin{tabularx}{\linewidth}{|>{\centering\arraybackslash}X|}
                  Grosor experimental \bolditalic{(mm)} \\\hline
                  $0,0750$                              \\\hline
            \end{tabularx}
      \end{center}
      \caption{}
      \label{tab:g}
\end{table}

% ! ----------------------------------------------------------------------|>
\section{Análisis de datos}

% + -------------------------------------------------------------|>
\subsection{Análisis}

% * ------------------------------------------------------|>
\subsubsection{}

\begin{table}[H]
      \begin{center}
            \begin{tabularx}{\linewidth}{|>{\centering\arraybackslash}X|>{\centering\arraybackslash}X|}
                  \hline
                  {\large $\lambda_{Teorico}$} \bolditalic{(nm)} & $632,8$ \\\hline
            \end{tabularx}
      \end{center}
\end{table}

\vspace{-.5cm}

\begin{table}[H]
      \begin{center}
            \begin{tabularx}{\linewidth}{|>{\centering\arraybackslash}X|>{\centering\arraybackslash}X|}
                  \multicolumn{2}{c}{$D_{1}$}          \\\hline
                  \textbf{n} & Error \bolditalic{(\%)} \\\hline
                  1          & $0,20$                  \\\hline

                  2          & $0,59$                  \\\hline

                  3          & $0,19$                  \\\hline

            \end{tabularx}
      \end{center}
      \caption{}
\end{table}

\begin{table}[H]
      \begin{center}
            \begin{tabularx}{\linewidth}{|>{\centering\arraybackslash}X|>{\centering\arraybackslash}X|}
                  \multicolumn{2}{c}{$D_{1}$}          \\\hline
                  \textbf{n} & Error \bolditalic{(\%)} \\\hline
                  1          & $1,86$                  \\\hline

                  2          & $0,51$                  \\\hline

                  3          & $0,53$                  \\\hline

            \end{tabularx}
      \end{center}
      \caption{}
\end{table}

% * ------------------------------------------------------|>
\subsubsection{}

Usando los datos de la tabla~(\ref{tab:g}) y la
formula~(\ref{eq:grosor-cabello}),

\begin{table}[H]
      \begin{center}
            \begin{tabularx}{\linewidth}{|>{\centering\arraybackslash}X|>{\centering\arraybackslash}X|}
                  \hline
                  Grosor Teorico \bolditalic{(mm)} & \multirow{2}{*}{$0,1171$} \\\hline
                  Error \bolditalic{(\%)}          & $35,9347$                 \\\hline
            \end{tabularx}
      \end{center}
\end{table}

% ! ----------------------------------------------------------------------|>
\section{Conclusiones}

\printbibliography

\end{document}